\section{Proof}
\begingroup
\footnotesize
\subsection{Direct Proof}

\textbf{1. If $n$ is odd, $n^2$ is odd.} \\
Assume that $n$ is odd. \\
Let $n=2k+1$, $k\in\mathbb Z$.
\[
n^2=(2k+1)^2=2(2k^2+2k)+1.
\]
Hence $n^2$ is odd. $\square$

\textbf{2. If $m,n$ are squares, $mn$ is a square.} \\
Assume that $m$ and $n$ are both squares. \\
Let $m=a^2,n=b^2$, $a,b\in\mathbb Z$.
\[
mn=a^2b^2=(ab)^2.
\]
Hence $mn$ is a perfect square. $\square$

\subsection{Proof by Contraposition}

\textbf{If $3n+2$ is odd, $n$ is odd.}

Contrapositive: \\ 
Assume that $n$ is even, $3n+2$ is even. \\
Let $n=2k$:
\[
3n+2=6k+2=2(3k+1).
\]
Since the negation of conclusion implies the negation of hypothesis, the original conditional statement is true. $\square$

\textbf{2. If $n=ab$, then $a\leq\sqrt n$
or $b\leq\sqrt n$.}

Contrapositive: \\ 
Assume that $a>\sqrt n$ and
$b>\sqrt n$, then $ab\neq n$.

Since $a,b>0$,
\[
ab>\sqrt n\cdot\sqrt n=n.
\]
Thus $ab\neq n$. \\
Since the negation of conclusion implies the negation of hypothesis, the original conditional statement is true. $\square$
\subsection{Proof by Contradiction}

\textbf{1. $\sqrt{2}$ is irrational.}

Assume that $\sqrt{2}$ is rational. \\
$\sqrt2=a/b$, 
$a,b\in\mathbb{Z},\quad b\neq 0,\quad a, b \text{ coprime }$
\[
a^2=2b^2\Rightarrow a=2k
\]
\[
4k^2=2b^2\Rightarrow b^2=2k^2.
\]
Both $a,b$ are even, they share a common factor of $2$. \\
It contradicts the initial assumption that $a, b$ \\ coprime, so the original statement is true.
$\square$

\textbf{2. If $3n+2$ is odd, $n$ is odd.}

Assume that the statement is false.
Then we have $3n+2$ is odd and $n$ is even.
Let $n=2k$:
\[
3n+2=2(3k+1),
\]
Hence $3n + 2$ is even. \\
It contradicts the initial assumption that $3n + 2$ is odd, so the original statement is true.
$\square$

\subsection{Proof of Equivalence}

\textbf{Prove:}
$P_1$: $n$ even;
$P_2$: $n-1$ odd;
$P_3$: $n^2$ even.

Show $P_1\Rightarrow P_2\Rightarrow P_3
\Rightarrow P_1$.

\textbf{$P_1\Rightarrow P_2$:}
$n=2k\Rightarrow n-1=2(k-1)+1$ (odd).

\textbf{$P_2\Rightarrow P_3$:}
$n-1=2k+1\Rightarrow n=2(k+1)$.
\[
n^2=2[2(k+1)^2]\quad\text{(even)}.
\]

\textbf{$P_3\Rightarrow P_1$:}
Contrapositive: $n=2k+1$ gives
\[
n^2=2(2k^2+2k)+1\quad\text{(odd)}.
\]
Hence all three are equivalent. $\square$

\subsection{Non-Constructive Proof}

\textbf{Exist irrational $x,y$ with rational $x^y$.}

Let $a=(\sqrt2)^{\sqrt2}$.

\textbf{Case 1:} If $a$ is rational, choose
$x=y=\sqrt2$. Then $x^y=(\sqrt2)^{\sqrt2}$ is rational.

\textbf{Case 2:} If $a$ is not rational, choose
$x=a$, $y=\sqrt2$:
\[
x^y=((\sqrt2)^{\sqrt2})^{\sqrt2}=2.
\], which is rational. \\
Both cases establish existence. $\square$

\subsection{Backward Reasoning}

\textbf{AM--GM:} For $a,b>0$, prove
$\frac{a+b}{2}\geq\sqrt{ab}$.

Work backward:
\[
\begin{aligned}
\frac{a+b}{2}\geq\sqrt{ab}
&\iff a+b\geq2\sqrt{ab}\\
&\iff a+b-2\sqrt{ab}\geq0\\
&\iff(\sqrt a-\sqrt b)^2\geq0.
\end{aligned}
\]
Last statement is always true.
Equality iff $a=b$. $\square$
\endgroup


